\textbf{CivilComments} (WILDS; \citealp{borkan2019nuanced,koh2021wilds}) pairs toxicity ($T$) with a human annotation that the comment mentions Muslims or LGBTQ people ($N$). Following \citet{dixon2018measuring} and \citet{garg2019counterfactual}, we swapped those identity terms for Christian/straight terms and, on 400 further comments, ``Christian'' for ``Muslim'', averaging the two directions; placebos also include ``Christian''$\to$``Protestant''. Both judges reject most of these swaps (Table~\ref{tab:data}: \civilAccFwd{} and \civilAccRev{} accepted by both), mostly because swapping the group changes a claim about it (``the baker discriminated against the gay couple''), not just its name. We therefore report Civil here only. Its results agree with the main settings: at $\rhot=0.9$ near-best extractors have $\cfn$ \nbCFmincivilQNine--\nbCFmaxcivilQNine{} (Qwen; Gemma \nbCFmincivilGNine--\nbCFmaxcivilGNine) and $\conf{}<0.5$ throughout, i.e.\ the most decodable toxicity scores rank non-toxic comments about Muslims or LGBTQ people above toxic comments that mention neither; at $\rhot=0.5$ the net shift is small (\nbCFmincivilQFive--\nbCFmaxcivilQFive); and choosing by \conf{} lowers $\cfn$ at a larger accuracy cost than in the main settings, because the nuisance is predictive in-distribution (Table~\ref{tab:near-all} and the selection table in Appendix~\ref{app:more}).

\begin{table}[h]
\centering\small\setlength{\tabcolsep}{2pt}
\resizebox{\columnwidth}{!}{\begin{tabular}{@{}llccccc@{}}
\toprule
 & & best & \#near & $\cfn$ near-best & $\cfn$ equiv. & \conf{} \\
 & & \dec{} & (equiv.) & range & range & range \\
\midrule
\multicolumn{7}{@{}l}{\emph{$\rhot=0.9$ (correlated training data)}}\\
CEBaB & Q & .966 & 37 (9) & .46--.89 & .54--.83 & .44--.69 \\
 & G & .979 & 39 (8) & .26--.72 & .25--.63 & .57--.82 \\
BiasBios & Q & .998 & 66 (65) & .09--.42 & .09--.42 & .83--.97 \\
 & G & .998 & 66 (61) & .09--.33 & .11--.33 & .92--.97 \\
Amazon-Neg & Q & .992 & 51 (45) & .03--.12 & .03--.12 & .87--.95 \\
 & G & .994 & 64 (60) & .02--.12 & .02--.11 & .87--.96 \\
Amazon-1P & Q & .993 & 40 (27) & .11--.25 & .11--.22 & .85--.95 \\
 & G & .994 & 62 (49) & .05--.31 & .06--.25 & .84--.96 \\
Civil & Q & .932 & 8 (5) & .24--.45 & .28--.42 & .35--.41 \\
 & G & .940 & 12 (1) & .30--.50 & .35 & .29--.43 \\
\midrule
\multicolumn{7}{@{}l}{\emph{$\rhot=0.5$ (decorrelated)}}\\
CEBaB & Q & .924 & 10 (1) & .02--.28 & .09 & .84--.90 \\
 & G & .960 & 11 (1) & $-$.10--.08 & .02 & .93--.95 \\
BiasBios & Q & .995 & 58 (54) & $-$.04--.05 & $-$.04--.05 & .98--.99 \\
 & G & .995 & 60 (53) & $-$.08--.03 & $-$.06--.03 & .98--.99 \\
Amazon-Neg & Q & .981 & 27 (14) & $-$.06--$-$.00 & $-$.05--$-$.01 & .95--.97 \\
 & G & .985 & 34 (23) & $-$.06--.00 & $-$.05--.00 & .95--.97 \\
Amazon-1P & Q & .986 & 28 (16) & $-$.01--.03 & $-$.01--.01 & .97--.98 \\
 & G & .987 & 53 (34) & $-$.01--.02 & $-$.01--.01 & .97--.99 \\
Civil & Q & .829 & 4 (2) & .04--.08 & $-$.00--.03 & .81--.83 \\
 & G & .839 & 5 (3) & $-$.02--.07 & $-$.03--.03 & .83--.84 \\
\bottomrule
\end{tabular}
}
\caption{As Table~\ref{tab:near}, all settings including Civil (accepted edits; Civil: mean of both swap directions).}
\label{tab:near-all}
\end{table}
